\documentclass[11pt]{article}

% --------------------------------------------------
% Geometry and Layout
% --------------------------------------------------
\usepackage[margin=1in]{geometry}
\usepackage{setspace}
\setstretch{1.15}

% --------------------------------------------------
% Mathematics and Symbols
% --------------------------------------------------
\usepackage{amsmath, amssymb, amsthm, mathtools}
\usepackage{bm}

% --------------------------------------------------
% Physics and Tensor Notation
% --------------------------------------------------
\usepackage{physics}

% --------------------------------------------------
% Fonts and Encoding
% --------------------------------------------------
\usepackage[T1]{fontenc}
\usepackage[utf8]{inputenc}

% --------------------------------------------------
% References and Links
% --------------------------------------------------
\usepackage{hyperref}
\hypersetup{
  colorlinks=true,
  linkcolor=blue,
  citecolor=blue,
  urlcolor=blue
}

% --------------------------------------------------
% Theorem Environments
% --------------------------------------------------
\theoremstyle{plain}
\newtheorem{theorem}{Theorem}
\newtheorem{lemma}{Lemma}
\theoremstyle{definition}
\newtheorem{definition}{Definition}

\title{Admissible Histories:\\
The Incongruent Neuron and the Ontological Structure of Error}
\author{Flyxion}
\date{\today}

\begin{document}
\maketitle

\section{Introduction: The Epistemic Shock of the Incongruent Neuron}

Recent work on biomimetic corticostriatal models reports the discovery of a class of neural signals whose activity predicts an erroneous behavioral outcome more than one second before that outcome is expressed. These so-called incongruent neurons fire reliably within the first two hundred milliseconds of stimulus onset, despite the fact that the associated motor response occurs an order of magnitude later in time. The epistemic shock of this finding lies not merely in its predictive power, but in the contradiction it poses to dominant accounts of error in neuroscience. In such accounts, error is typically treated as noise, stochastic drift, or late-stage failure to converge on a target representation. Under that view, errors arise when computation breaks down.

The existence of early, stable, causally efficacious error-predictive activity renders such failure-centric models untenable. The incongruent neuron does not mark the absence of computation; it marks the presence of a fully specified internal process that unfolds coherently toward an outcome that is externally evaluated as incorrect. The central thesis of this paper is therefore ontological rather than descriptive: the incongruent neuron authorizes an admissible history. What appears as error at the level of behavior is, at the level of neural dynamics, the successful execution of a complete and internally coherent trajectory that happens to be disfavored by task contingencies.

\subsection*{Formal Derivation}

Let a trial be modeled as a finite sequence of neural states
\[
h = (x_0, x_1, \dots, x_T),
\]
with $x_0$ corresponding to stimulus onset and $x_T$ to behavioral expression. Define a mapping
\[
\pi : h \mapsto a \in \mathcal{A},
\]
from histories to actions. The empirical finding implies that for some $t \ll T$, the conditional probability
\[
\mathbb{P}(\pi(h)=a_{\mathrm{err}} \mid x_0,\dots,x_t) \approx 1.
\]
Thus, the outcome is fixed by a prefix of the history. Error is not a deviation at $T$, but a property of the prefix $(x_0,\dots,x_t)$ itself.

\section{The Temporal Collapse: Prefix-Closure and Commitment}

The most striking temporal feature of the incongruent neuron is the collapse of decision time. Within approximately two hundred milliseconds, the system enters a dynamical regime from which the eventual outcome can be inferred with high confidence. Behavioral response latency therefore ceases to be an indicator of decision formation and becomes instead a lagging indicator of a commitment already made.

This phenomenon is naturally described by the notion of prefix-closure. A set of histories is prefix-closed if, whenever a history is admissible, all of its initial segments are likewise admissible. In the present context, the early neural activity constitutes a prefix that constrains all future continuations. Once this prefix is realized, the space of possible futures collapses onto a narrow subset consistent with that initial commitment. What is experienced phenomenologically as deliberation is, under this view, the deterministic unfolding of an already authorized trajectory.

\subsection*{Formal Derivation}

\begin{definition}
Let $\mathcal{H}$ be a set of finite histories. We say $\mathcal{H}$ is prefix-closed if for any $h = (x_0,\dots,x_T)\in\mathcal{H}$ and any $t<T$, the prefix $h_{\le t}=(x_0,\dots,x_t)$ is also in $\mathcal{H}$.
\end{definition}

\begin{lemma}
If a history $h$ is uniquely determined by its prefix $h_{\le t}$, then no extension of $h_{\le t}$ outside the equivalence class of $h$ is admissible.
\end{lemma}

\begin{proof}
Uniqueness implies that for any extension $h'$ of $h_{\le t}$, $h' \neq h$ violates the system dynamics. Hence $h'\notin\mathcal{H}$.
\end{proof}

The incongruent neuron identifies precisely such prefixes.

\section{From Entities to Trajectories: The Logic of Admissibility}

Incongruent neurons do not encode error as an abstract category. They encode alternative action paths with the same structural fidelity as neurons associated with correct outcomes. The distinction between congruent and incongruent activity is therefore not representational but evaluative: it arises only when histories are assessed relative to external reward signals.

This motivates a formal distinction between coherence and value. Coherence is an internal property, defined by whether a trajectory is dynamically permitted by the architecture. Value is an external assignment, imposed by feedback. Admissibility concerns the former, not the latter. The space of admissible histories $\mathcal{H}$ is defined by the constraints of the corticostriatal loop, including lateral inhibition and striatal gating. Incorrect histories belong to $\mathcal{H}$ no less than correct ones; they differ only in valuation.

\subsection*{Formal Derivation}

\begin{definition}
Let $\mathcal{H}$ denote the set of all histories consistent with neural dynamics. Let $V:\mathcal{H}\to\mathbb{R}$ be a valuation function induced by reward.
\end{definition}

\begin{theorem}
Membership in $\mathcal{H}$ is independent of the sign of $V(h)$.
\end{theorem}

\begin{proof}
By construction, $\mathcal{H}$ is determined solely by dynamical constraints. The valuation $V$ is applied after the fact and does not affect admissibility.
\end{proof}

Thus, incorrectness does not negate admissibility.

\section{Structural Inevitability: Why Incongruent Neurons Are Mandated}

Given sparse cortical connectivity and biased receptive fields, initial conditions are unavoidably heterogeneous. Lateral inhibition amplifies these small differences into discrete competitive outcomes. Reinforcement learning then sharpens synaptic weights without erasing alternative paths. As a result, losing trajectories are preserved as latent but executable histories.

Under these conditions, neurons that preferentially support unrewarded outcomes are not anomalies but mathematical necessities. Any architecture that combines competition, learning, and irreversibility must generate specialized carriers of disfavored trajectories. The incongruent neuron is the explicit manifestation of this requirement.

\subsection*{Formal Derivation}

Let synaptic weights evolve as
\[
w_{ij}(t+1)=w_{ij}(t)+\eta r(t)\Delta_{ij}(t),
\]
with learning rate $\eta$ and reward $r(t)$. If $\Delta_{ij}(t)\neq 0$ for both correct and incorrect paths, then weights diverge but do not vanish. Hence,
\[
\exists h_{\mathrm{err}}\in\mathcal{H}\quad\text{with}\quad V(h_{\mathrm{err}})<0.
\]

\section{The Intervention Proof: Pruning the Future}

The strongest empirical evidence for the ontological status of error histories comes from intervention. When trials are halted in real time upon detection of incongruent neuron activity, behavioral accuracy increases significantly. This intervention does not correct an unfolding error; it removes an entire future from realization. The operation is therefore best understood as an exclusion operator acting on $\mathcal{H}$.

By preventing the completion of a history whose prefix has already been decoded, the intervention demonstrates that the error was already fully present in the neural substrate. Accuracy improves not by repairing computation, but by pruning the space of executable histories.

\subsection*{Formal Derivation}

Let $\mathcal{H}'\subset\mathcal{H}$ be the set remaining after intervention. Accuracy is
\[
\mathrm{Acc}=\frac{|\{h\in\mathcal{H}':V(h)>0\}|}{|\mathcal{H}'|}.
\]
Removing histories with $V(h)<0$ strictly increases $\mathrm{Acc}$ without altering dynamics on $\mathcal{H}'$.

\section{Conclusion: Error as a Legitimate History}

The cumulative evidence forces a conceptual shift. The brain does not make mistakes in the sense of computational failure. It executes histories. Some of these histories are rewarded; others are not. Rationality and correctness emerge as algebraic properties of successful pruning, not as axioms of neural computation.

The incongruent neuron reveals the primary unit of cognition to be the authorized trajectory. A neuroscience of error that ignores this fact studies shadows rather than structure. The proper object of inquiry is the mechanism by which histories are authorized, weighted, and excluded before they reach behavioral expression.

\end{document}